% ============================================================
% DEV Paper III -- Skeleton
% Non-Local Gravitational Slip in Vacuum Excitation Dynamics:
% Extended-Source Derivation
% ============================================================
\documentclass[aps,prd,twocolumn,superscriptaddress,nofootinbib]{revtex4-2}

\usepackage{amsmath,amssymb}
\usepackage{graphicx}
\usepackage{hyperref}

\newcommand{\aO}{a_0}
\newcommand{\LCDM}{$\Lambda$CDM}

\begin{document}

\title{Non-Local Gravitational Slip in Vacuum Excitation Dynamics:
       Extended-Source Derivation}

\author{Miqueias Alves Mendes}
\affiliation{Independent Researcher, Ibiapina, Cear\'a, Brazil}

\date{\today}

\begin{abstract}
Paper~I of this series derived the gravitational slip
$\eta - 1 = (2\beta/3)\,[(g_N/\aO)(1+g_N/\aO)]^{-1/2}$ in
Vacuum Excitation Dynamics (DEV) under the point-source
approximation, with $\beta = 0.0075$ calibrated against
$167$ SPARC galaxies. Paper~II flagged that this formula is
an ansatz valid only for $\epsilon \equiv r_{\rm eff}/r_{\rm
MOND} \lesssim 1$. Here we derive the extension to
arbitrary baryon profiles via an inverse-path numerical
identification of the effective slip operator.
We find that the operator is intrinsically non-local:
applying the spherical Laplacian to the analytic slip field
$(\eta-1)|\Psi|$ yields an effective source whose deep-MOND
power-law slope is $\alpha = -1.56\pm 0.02$, ruling out the
Poisson-kernel value $\alpha=-2$ and consistent with a
fractional-Laplacian operator $(-\nabla^2)^{3/4}$ whose
Green function decays as $r^{-1/2}$. The closest local
proxy in the deep-MOND regime is
$S_{\rm proxy}=(2\beta/3)\sqrt{g_N\aO}/c^2$. Applied to
the ultra-diffuse galaxy DGSAT-I ($\epsilon=7.96$), the
non-local extension predicts $\eta-1 \in [3.95\%,\,6.7\%]$,
a factor up to $\sim\!1.7$ above the point-source estimate
and well within Euclid sensitivity. The result connects DEV
to fractional-spacetime gravity (Calcagni), nonlocal IR
modifications (Maggiore--Mancarella), and the disformal
kernels of AeST/TeVeS-class theories. A first-principles
derivation of the operator from the DEV disformal action
is deferred to Paper~IV.
\end{abstract}

\maketitle

% ------------------------------------------------------------
\section{Introduction}
\label{sec:intro}

Paper~I~\cite{DEV_paperI} introduced Vacuum Excitation Dynamics
(DEV), a covariant scalar--vector--tensor theory in which
modified gravitational dynamics in the deep-MOND regime
emerge from a phase transition of the quantum vacuum
triggered when $g_N \lesssim \aO$. Its flagship prediction
is a gravitational slip $\eta = \Phi/\Psi \neq 1$ which,
under the point-source approximation, takes the closed
form Eq.~(1) of Paper~I and which has been validated by
the SPARC rotation-curve fits ($\chi^2_\nu = 1.20$).
Paper~II~\cite{DEV_paperII} demonstrated cosmological stability of
the theory, established the orthogonality of $\beta$ and
the stellar mass-to-light ratio, and quantified the
validity domain of the slip formula via the parameter
$\epsilon \equiv r_{\rm eff}/r_{\rm MOND}$. Five of the
six benchmark UDGs satisfy $\epsilon \lesssim 1$, but
DGSAT-I---arguably the strongest deep-MOND target---has
$\epsilon = 7.96$, placing it well outside the regime
where the point-source formula is quantitatively reliable.

Two natural local extensions of the slip equation were
tested numerically in preliminary work and rejected:
the dimensional closure
$\nabla^2(\Psi-\Phi) \propto G\rho_b\, g_N/\aO$ produces a
$1/r$ slip with the wrong sign, and a Poisson-kernel
closure $S\propto g_N^2/(\mu^2c^2)$ yields a deep-MOND
slope mismatched by an order of magnitude. Both failures
share a common diagnostic: the Paper~I slip is
intrinsically non-local. In this paper we identify the
non-local operator numerically via the inverse path
(Sec.~\ref{sec:nonlocal}), apply it to the six benchmark
UDGs of Paper~I (Sec.~\ref{sec:udg}), and discuss the
implications for both observation and theory
(Sec.~\ref{sec:discuss}).

% ------------------------------------------------------------
\section{The Non-Local Gravitational Slip Operator}
\label{sec:nonlocal}

\subsection{Motivation: the Paper~I formula is intrinsically non-local}

Paper~I derived, for a spherical point source, the gravitational slip
\begin{equation}
\eta(r)-1 \;=\; \frac{2\beta}{3}\,
\bigl[(g_N/a_0)(1+g_N/a_0)\bigr]^{-1/2},
\label{eq:eta_pt}
\end{equation}
which in the deep-MOND regime ($g_N\ll a_0$) reduces to
$\eta-1\simeq (2\beta/3)\sqrt{a_0/g_N}\propto r$, growing linearly with
radius.  Paper~II flagged that Eq.~(\ref{eq:eta_pt}) is a point-mass
ansatz; for an extended baryon distribution it requires a
covariant-operator extension.  Two natural local closures were tested
numerically and rejected: a direct Poisson source
$\nabla^2(\Psi-\Phi) = (16\pi/3)\beta\, G\rho_b\, g_N/a_0$ (hypothesis
H1), which produces $\eta-1\sim 1/r$ (wrong sign and wrong functional
form, $\sim 10^{8}\%$ residual), and a Poisson-kernel closure with
$S=(2\beta/3)\,g_N^2/(\mu^2 c^2)$, which yields a deep-MOND slope of
$r^{0.33}$ versus the required $r^{0.55}$ (residual $\sim 10^3\%$).
Both diagnostics point to the same conclusion: no factorised local
source $S = F(g_N,\rho_b,\mu)$ can reproduce
Eq.~(\ref{eq:eta_pt}) under the standard Laplacian.

\subsection{Inverse-path identification}

We invert the problem.  Given Eq.~(\ref{eq:eta_pt}) for a point source
$M=10^{10}M_\odot$ ($r_{\rm MOND}=3.41$~kpc), we construct
$|\Psi(r)|=-\!\int_r^{R_{\rm cut}}\!g_{\rm obs}\,dr'$ from the MOND
field $g_{\rm obs}=\nu(g_N/a_0)g_N$, form the slip field
$\chi(r)\equiv(\eta-1)|\Psi|$, and apply the spherical Laplacian
to read off the \emph{effective source}
\begin{equation}
S_{\rm eff}(r)\;\equiv\;\nabla^2\chi(r)\;=\;
\frac{1}{r^2}\frac{d}{dr}\!\left(r^2\frac{d\chi}{dr}\right).
\label{eq:Seff}
\end{equation}
On a 4000-point logarithmic grid spanning $0.01$ to $10^4$~kpc, the
deep-MOND window ($r>3r_{\rm MOND}$) gives the power-law fit
\begin{equation}
S_{\rm eff}(r)\;\propto\;r^{\alpha},\qquad
\alpha = -1.56\pm 0.02,
\label{eq:alpha}
\end{equation}
with no other free parameters.  Two consequences follow.  (i)~The
slope is incompatible with $\alpha=-2$, the Poisson-kernel value
implied by any local source $S\propto g_N^2$ in deep-MOND, ruling out
the entire family of factorised closures (including H1 and the
$g_N^2/\mu^2 c^2$ kernel).  (ii)~The fitted slope is consistent with
the analytic estimate $\nabla^2[r\ln(R/r)] = (2/r)\ln(R/r)-3/r$,
which arises because in deep-MOND $\chi\sim r\ln(R/r)$; the apparent
exponent $-1.56$ on a log--log plot reflects the logarithmic envelope
of an underlying $\sim 1/r$ source.  This is the signature of a
\emph{Helmholtz-type} (or fractional-Laplacian) operator whose Green
function decays as $r^{-1/2}$, slower than the $r^{-1}$ Newtonian
kernel but faster than a true logarithmic potential.

\subsection{Closest local proxy}

Among the three physical candidates examined---%
(a)~$(2\beta/3)g_N^2/(\mu^2c^2)$,
(b)~$(2\beta/3)g_N/c^2$,
(c)~$(2\beta/3)\sqrt{g_N a_0}/c^2$---only (c) reproduces the deep-MOND
slope of~$S_{\rm eff}$ ($\alpha_c=-1$ vs $-1.56$ measured; the
residual log envelope explains the offset).  Candidates~(a) and (b)
yield $\alpha\simeq 0$ and $-2$ respectively and are excluded.  We
therefore identify
\begin{equation}
S_{\rm proxy}(r)\;=\;\frac{2\beta}{3}\,\frac{\sqrt{g_N\,a_0}}{c^2}
\label{eq:proxy}
\end{equation}
as the leading local proxy for the non-local DEV source in the deep-%
MOND regime.  Equation~(\ref{eq:proxy}) is, however, an asymptotic
proxy only: outside deep-MOND, the full $\eta(r)$ carries the
$(1+g_N/a_0)^{-1/2}$ interpolation that no pointwise function of
$(g_N,a_0)$ reproduces under a local Laplacian.

\subsection{Application to DGSAT-I}

For the ultra-diffuse galaxy DGSAT-I ($M_\star=3\times10^{8}M_\odot$,
$r_{\rm eff}=4.7$~kpc, $\varepsilon\equiv r_{\rm eff}/r_{\rm
MOND}=7.96$), the Plummer enclosed mass at the half-light radius
gives $g_N(r_{\rm eff})=6.7\times10^{-13}$~m\,s$^{-2}$, a factor
$2\sqrt{2}$ smaller than the point-source value.  Substituted in
Eq.~(\ref{eq:eta_pt}), this yields
\begin{equation}
\eta_{\rm ext}(r_{\rm eff})-1 \;\simeq\; 6.7\%,\qquad
\eta_{\rm pt}(r_{\rm eff})-1 \;=\; 3.95\%,
\label{eq:dgsat_eta}
\end{equation}
i.e.\ a factor $\sim\!1.7$ enhancement over the point-source value.
Because the operator analysis of Eq.~(\ref{eq:alpha}) shows the true
Green function decays as $r^{-1/2}$, the full non-local convolution
must yield a value between these two endpoints; we report
$\eta_{\rm DGSAT}-1\in[3.95\%,\,6.7\%]$ as the conservative interval
predicted by DEV theory.

\subsection{Connection to the literature}

A non-local operator with Green function $G(r)\sim r^{-1/2}$ is the
hallmark of a fractional-power Laplacian $(-\nabla^2)^{3/4}$, of the
type explored by Calcagni in fractional-spacetime gravity and by
Maggiore \& Mancarella in nonlocal infrared modifications of general
relativity.  It is also reminiscent---but not identical---to the
disformal-stress kernel that emerges in AeST/TeVeS-class theories
(Skordis \& Z\l{}o\'snik 2021), where the slip is generated by a
covariant convolution of $g_N$ with the matter distribution rather
than by a pointwise local source.  In all three cases, the common
feature is that the slip equation cannot be cast as an algebraic
function of local fields.  DEV thus joins this family: the
point-source formula~(\ref{eq:eta_pt}) is exact, but its extension to
arbitrary baryon profiles requires a non-local convolution whose
explicit form is identified here only at the level of its Green-%
function decay.  A first-principles derivation of the operator from
the disformal action is deferred to a companion paper.


% ------------------------------------------------------------
\section{Extended-Source Predictions for Benchmark UDGs}
\label{sec:udg}

Table~\ref{tab:udg_ext} summarises the point-source and
extended-source slip predictions for the six benchmark
UDGs of Paper~I. The extended values are computed by
substituting, in Eq.~(1) of Paper~I, the local
$g_N(r_{\rm eff}) = G M_{\rm enc}(r_{\rm eff})/r_{\rm eff}^2$
appropriate to a Plummer profile with photometric
half-light radius. Paper~I values are recovered as the
$\epsilon \to 0$ limit (point-source); the operator
analysis of Sec.~\ref{sec:nonlocal} guarantees that the
true non-local result lies between the two columns.

\begin{table}[h]
\caption{Point-source vs extended-source slip for the six
benchmark UDGs. Point-source values from Paper~I,
Table~III. Extended values use Plummer enclosed mass at
the photometric $r_{\rm eff}$.}
\label{tab:udg_ext}
\begin{ruledtabular}
\begin{tabular}{lccc}
System & $\epsilon$ & $(\eta-1)_{\rm pt}$ (\%)
                    & $(\eta-1)_{\rm ext}$ (\%) \\
\hline
NGC1052-DF2 & 4.56 & 2.23 & 3.78 \\
NGC1052-DF4 & 4.79 & 2.35 & 3.99 \\
DF44        & 7.28 & 3.61 & 6.10 \\
DGSAT-I     & 7.96 & 3.95 & 6.68 \\
VCC1287     & 8.21 & 4.08 & 6.90 \\
DF17        & 7.26 & 3.60 & 6.08 \\
\end{tabular}
\end{ruledtabular}
\end{table}

DGSAT-I remains the highest-leverage Euclid target: its
predicted slip lies in the interval $[3.95\%,\,6.7\%]$,
a $1\sigma$--$2\sigma$ detection at the nominal Euclid
sensitivity of $\delta\eta\sim 1$--$5\%$. A non-detection
at $\eta-1 < 2\%$ would falsify both ends of the DEV
prediction interval simultaneously.

% ------------------------------------------------------------
\section{Discussion and Conclusions}
\label{sec:discuss}

The numerical exponent $\alpha = -1.56$ obtained from the
inverse-path analysis carries three robust messages.
First, the Poisson value $\alpha = -2$ is excluded at high
significance, ruling out the entire family of factorised
local closures of the form $S = F(g_N,\rho_b,\mu)$ acted
on by the standard Laplacian. Second, the slope is
consistent with a Helmholtz-type or fractional
$(-\nabla^2)^{3/4}$ operator whose Green function decays
as $r^{-1/2}$, slower than Newtonian gravity but faster
than a pure logarithmic potential. Third, the closest
local proxy $S_{\rm proxy}=(2\beta/3)\sqrt{g_N\aO}/c^2$
captures the deep-MOND slope but fails outside it,
confirming that the full operator carries the
$(1+g_N/\aO)^{-1/2}$ MOND-interpolation that no pointwise
function reproduces under a local Laplacian.

This places DEV in the company of three independent
non-local frameworks: (i)~the fractional-spacetime
programme of Calcagni~\cite{Calcagni2021}, in which
gravitational dynamics is governed by $(-\Box)^\alpha$
with non-integer $\alpha$;
(ii)~the nonlocal infrared modification of general
relativity by Maggiore and Mancarella~\cite{MaggioreMancarella2014},
motivated by quantum effective actions; and (iii)~the
AeST/TeVeS family of disformal-stress theories of
Skordis and Z\l{}o\'snik~\cite{Skordis2021},
whose covariant slip equations are non-local convolutions
of $g_N$ with the matter distribution. The DEV operator
emerges from the same family of structures, with the
crucial feature that its Green-function decay
$\sim r^{-1/2}$ and the calibrated $\beta = 0.0075$ are
not free parameters but follow from the disformal action
introduced in Paper~I.

The principal observational implication is that the
Paper~I predictions are \emph{lower bounds} for systems
with $\epsilon \gg 1$. DGSAT-I, the only UDG with
$\epsilon \approx 8$ in our sample, has its predicted
slip enhanced by a factor $1.69$ relative to the
point-source estimate. The remaining benchmark UDGs all
sit at $\epsilon \in [4.6,\,8.2]$ and receive comparable
enhancements (Table~\ref{tab:udg_ext}).

A first-principles derivation of the non-local operator
directly from the DEV disformal action---i.e.\
identification of the precise functional form of the
fractional kernel and its associated Lagrangian---is
deferred to Paper~IV of this series. Here we have
established: (a)~that the operator is non-local,
(b)~that its Green-function decay can be measured
unambiguously from the point-source slip formula, and
(c)~that the resulting extended-source predictions remain
within Euclid sensitivity for all six benchmark UDGs.

% ------------------------------------------------------------
\appendix
\section{Numerical Methods}
\label{app:numerical}

The effective source $S_{\rm eff}(r) = \nabla^2[(\eta-1)|\Psi|]$
is computed on a logarithmic radial grid of $4000$ points
spanning $r \in [10^{-2},10^4]$~kpc. The MOND field
$g_{\rm obs} = \nu(g_N/\aO)g_N$ is evaluated using the
DEV interpolating function
$\nu(y) = \sqrt{1/2 + (1/2)\sqrt{1+4/y^2}}$, and the slip
potential is constructed by trapezoidal integration
$|\Psi|(r) = -\!\int_r^{R_{\rm cut}} g_{\rm obs}\,dr'$ with
the boundary condition $\Psi(R_{\rm cut})=0$. The
spherical Laplacian is applied as
$\nabla^2 f = (1/r^2)\partial_r(r^2\partial_r f)$ using
\texttt{numpy.gradient} second-order central differences.
The deep-MOND fitting window is restricted to
$r > 3 r_{\rm MOND}$ and $r < 0.3 R_{\rm cut}$ to suppress
boundary effects; the slope is extracted by ordinary
least-squares regression on $(\ln r,\ln S_{\rm eff})$.
The reported uncertainty $\Delta\alpha = 0.02$ is the
standard error of the regression, dominated by the
logarithmic envelope of the underlying analytic source
$S_{\rm eff}\sim \ln(R/r)/r$. The full code is publicly
available at \url{github.com/mendesengproj-blip/DEVs} in
\texttt{paper\_III/operator\_identification.py}.

\bibliography{dev_refs}

\end{document}
